\section{Roteiro Prático 04 — Arquitetura Frontend Modular: ES6 Modules \& Padrão em Camadas}

% =====================================================
% OBJETIVOS DA ATIVIDADE
% =====================================================

\begin{boxObjetivos}
\textbf{Objetivos da Atividade (Módulo 2 — ES6 Modules \& Padrão em Camadas)}

Ao final desta prática orientada em laboratório, o estudante será capaz de:

\begin{itemize}
\item Superar a arquitetura monolítica de arquivo único (\texttt{app.js}), organizando o frontend em múltiplos módulos especializados utilizando \textbf{ES6 Modules} (\texttt{import} e \texttt{export});
\item Declarar o ponto de entrada modular na página web através da diretiva:\\ \texttt{<script type="module" src="js/main.js"></script>};
\item Compreender e aplicar a separação estrita de responsabilidades no cliente em três camadas:
  \begin{itemize}
    \item \textbf{Camada de Serviços (\texttt{services/}):} Regras de negócio, cálculos e persistência pura no \texttt{localStorage} sem acoplamento ao DOM;
    \item \textbf{Camada de Visão (\texttt{views/}):} Renderização de cards, tabelas, modais e Toasts via \textit{template literals} do Bootstrap 5;
    \item \textbf{Camada Utilitária (\texttt{utils/}):} Funções puras reutilizáveis de formatação (moeda brasileira, telefone, máscaras);
    \item \textbf{Orquestrador Central (\texttt{main.js}):} Captura de eventos do usuário e coordenação entre camadas.
  \end{itemize}
\item Implementar o \textbf{CRUD Completo} (\textit{Create, Read, Update, Delete}) com persistência em armazenamento local;
\item Implementar \textbf{validação visual no cliente} com feedback imediato ao usuário utilizando as classes nativas do Bootstrap 5 (\texttt{.is-valid}, \texttt{.is-invalid} e \texttt{.invalid-feedback});
\item Disparar notificações flutuantes dinâmicas (\textbf{Bootstrap Toasts}) para confirmar operações de cadastro, alteração e exclusão;
\item Implementar filtragem reativa em tempo real por categoria de serviço na vitrine comunitária.
\end{itemize}
\end{boxObjetivos}

% =====================================================
% INTRODUÇÃO E CONTEXTUALIZAÇÃO
% =====================================================

\subsection{Introdução e Contextualização do Projeto}

Nos Roteiros 1, 2 e 3, trabalhamos com scripts JavaScript concentrados em um único arquivo (\texttt{app.js}). Embora essa abordagem seja excelente para aprender os fundamentos, ela se torna inviável em aplicações reais de médio e grande porte, onde a mistura de seletores do DOM, regras de validação e persistência gera o chamado \textit{"código espaguete"}.

No \textbf{Roteiro Prático 04}, damos o passo definitivo rumo à arquitetura profissional de software no frontend. Desenvolveremos a \textbf{Vitrine Comunitária de Empreendedores Locais de Garopaba}, uma aplicação que permite cadastrar, listar, filtrar, editar e gerenciar serviços da comunidade local. 

Este projeto serve também como a \textbf{base tecnológica inicial para a Ação Curricularizada de Extensão na SNCT (Semana Nacional de Ciência e Tecnologia)}.

% =====================================================
% ESTRUTURA DE DIRETÓRIOS DO PROJETO
% =====================================================

\subsection{Estrutura Arquitetural de Pastas}

O projeto deve ser organizado na pasta \texttt{roteiro-4/exercicio-roteiro4/} seguindo rigorosamente a estrutura modular:

\begin{boxCodigo}
\textbf{Árvore de Diretórios Modular do Projeto:}

\footnotesize
\begin{verbatim}
roteiro-4/exercicio-roteiro4/
|-- index.html     <-- Vitrine e forms (BS5)
|-- styles.css     <-- Estilos e transições
`-- js/
    |-- main.js    <-- Orquestrador central
    |-- services/
    |   `-- vitrineService.js <-- CRUD Storage
    |-- views/
    |   `-- vitrineView.js    <-- HTML dinâmico
    `-- utils/
        `-- formatters.js     <-- Formatações
\end{verbatim}
\end{boxCodigo}

% =====================================================
% FUNDAMENTAÇÃO TEÓRICA: ES6 MODULES
% =====================================================

\subsection{Fundamentação Teórica: Como Funcionam os ES6 Modules}

\subsubsection{1. Declaração no HTML (\texttt{type="module"})}
Ao declarar \texttt{<script type="module" src="js/main.js"></script>}, o navegador aplica comportamentos modernos essenciais:
\begin{itemize}
  \item \textbf{Modo Estrito Automático (\texttt{"use strict"}):} Impede o uso de variáveis não declaradas e falhas silenciosas;
  \item \textbf{Escopo Fechado de Módulo:} Variáveis e funções criadas dentro de um arquivo não poluem o escopo global (\texttt{window});
  \item \textbf{Carregamento Assíncrono com Adiamento (\texttt{defer}):} O script aguarda o parsing completo do DOM antes de executar.
\end{itemize}

\subsubsection{2. Exportação e Importação (\texttt{export} e \texttt{import})}
Para compartilhar recursos entre arquivos, usamos a sintaxe explícita de módulos:

\textbf{Exemplo de Exportação Nomeada e Importação:}

\begin{boxCodigoJS}
// No arquivo js/utils/formatters.js
export function formatarTelefone(telefone) {
  return telefone.replace(
    /^(\d{2})(\d{5})(\d{4})/,
    '($1) $2-$3'
  );
}

// No arquivo js/main.js
import { formatarTelefone }
  from './utils/formatters.js';

// Resultado: (48) 99999-8888
console.log(formatarTelefone('48999998888'));
\end{boxCodigoJS}

% =====================================================
% PIPELINE VISUAL DA ATIVIDADE
% =====================================================

\providecolor{ifscGreen}{RGB}{22,163,74}
\vspace{0.4cm}
\begin{figure}[H]
\centering
\resizebox{\linewidth}{!}{%
\begin{tikzpicture}[
    node distance=0.8cm and 0.5cm,
    every node/.style={font=\small},
    etapa/.style={
        rectangle,
        rounded corners=5pt,
        draw=azulTitulo,
        thick,
        minimum width=4.9cm,
        minimum height=1.75cm,
        text width=4.7cm,
        align=center,
        fill=azulTitulo!6
    },
    etapaHighlight/.style={
        rectangle,
        rounded corners=5pt,
        draw=ifscGreen,
        thick,
        minimum width=4.9cm,
        minimum height=1.75cm,
        text width=4.7cm,
        align=center,
        fill=ifscGreen!10
    },
    seta/.style={
        -{Latex[length=2.5mm]},
        thick,
        draw=azulTitulo
    }
]

% Linha 1: Estrutura Base e Serviços
\node[etapa] (p1) {\textbf{Passo 1: HTML5 \& Bootstrap}\\\footnotesize Estrutura base, modais, toasts e\\\texttt{<script type="module">}};

\node[etapa, right=of p1] (p2) {\textbf{Passo 2: Utilitários Puros}\\\footnotesize \texttt{formatters.js}: moedas, fone e\\\texttt{escapeHtml} contra XSS};

\node[etapa, right=of p2] (p3) {\textbf{Passo 3: Camada de Serviços}\\\footnotesize \texttt{vitrineService.js}: CRUD puro\\(\textit{Create, Read, Update, Delete})};

% Linha 2: Interface e Orquestração
\node[etapa, below=of p3] (p4) {\textbf{Passo 4: Camada de Visão}\\\footnotesize \texttt{vitrineView.js}: cards dinâmicos,\\estado vazio e feedback Toast};

\node[etapa, left=of p4] (p5) {\textbf{Passo 5: Orquestrador Central}\\\footnotesize \texttt{main.js}: validação nativa BS5,\\filtros reativos e eventos CRUD};

\node[etapaHighlight, left=of p5] (p6) {\textbf{Conclusão \& Deploy}\\\footnotesize Teste do ciclo completo e deploy\\automático no \textbf{GitHub Pages}};

% Conexões do fluxo
\draw[seta] (p1) -- (p2);
\draw[seta] (p2) -- (p3);
\draw[seta] (p3) -- (p4);
\draw[seta] (p4) -- (p5);
\draw[seta] (p5) -- (p6);

\end{tikzpicture}%
}
\caption{Pipeline arquitetural e sequência cronológica de implementação do Roteiro 04.}
\end{figure}

% =====================================================
% PASSO A PASSO GUIADO DA ATIVIDADE
% =====================================================

\clearpage
\subsection{Passo a Passo Guiado do Roteiro}

\begin{boxAtividade}[{Fase 1.1/5: Estrutura HTML5 Base e Ponto de Entrada (index.html)}]
\textbf{Responsabilidade Arquitetural:} Estabelecer o esqueleto visual da página, carregar as dependências de estilo (Bootstrap 5 e Icons), fornecer os pontos de ancoragem da interface (\texttt{\#vitrineContainer}, \texttt{\#toastNotificacao}, \texttt{\#totalServicosBadge} e \texttt{\#filtroCategoria}) e conectar o módulo raiz \texttt{js/main.js}.
\end{boxAtividade}

\vspace{-0.1cm}
\begin{boxCroqui}[{Croqui Visual 1.1: Wireframe da Interface no Navegador}]
+-----------------------------------------------------------------------+
| [Vitrine de Garopaba] (0 cadastrados)             [+ Divulgar Servico] | <- Header
+-----------------------------------------------------------------------+
| [ Filtro: Todas as Categorias v ]                 [ Toast Notificacao] | <- Controles
+-----------------------------------------------------------------------+
|                                                                       |
|    <div id="vitrineContainer">: Grid vazio aguardando os cards...     | <- Grid
|                                                                       |
+-----------------------------------------------------------------------+
\end{boxCroqui}

\begin{boxCodigoHTML}
<!-- index.html (Parte 1: Estrutura Base e Scripts) -->
<!DOCTYPE html>
<html lang="pt-BR">
<head>
  <meta charset="UTF-8">
  <meta name="viewport" content="width=device-width, initial-scale=1.0">
  <title>Vitrine Comunitária de Garopaba — Roteiro 04</title>
  <!-- Bootstrap 5 CSS & Icons CDN -->
  <link rel="stylesheet"
    href="https://cdn.jsdelivr.net/npm/bootstrap@5.3.3/dist/css/bootstrap.min.css">
  <link rel="stylesheet"
    href="https://cdn.jsdelivr.net/npm/bootstrap-icons@1.11.3/font/bootstrap-icons.min.css">
  <link rel="stylesheet" href="styles.css">
</head>
<body class="bg-light">

  <!-- Container de Toasts com corpo #toastMensagem -->
  <div class="toast-container position-fixed top-0 end-0 p-3" style="z-index: 1090;">
    <div id="toastNotificacao"
         class="toast align-items-center text-bg-success border-0 shadow"
         role="alert">
      <div class="d-flex">
        <div class="toast-body d-flex align-items-center gap-2">
          <i class="bi bi-check-circle-fill"></i>
          <div>
            <strong id="toastTitulo" class="d-block">Sucesso!</strong>
            <span id="toastMensagem">Operação realizada com êxito.</span>
          </div>
        </div>
        <button type="button" class="btn-close btn-close-white me-2 m-auto"
                data-bs-dismiss="toast"></button>
      </div>
    </div>
  </div>

  <!-- Header / Barra Superior com Contador e Botão de Ação -->
  <header class="bg-success text-white py-3 mb-4 shadow-sm">
    <div class="container d-flex justify-content-between align-items-center flex-wrap gap-2">
      <div>
        <h1 class="h4 mb-0 fw-bold">Vitrine de Garopaba</h1>
        <small class="opacity-75">Serviços e comércios locais</small>
      </div>
      <div class="d-flex align-items-center gap-2">
        <span class="badge bg-light text-dark" id="totalServicosBadge">0 cadastrados</span>
        <button class="btn btn-warning fw-bold" data-bs-toggle="modal"
                data-bs-target="#modalCadastro">
          <i class="bi bi-plus-circle me-1"></i> Divulgar
        </button>
      </div>
    </div>
  </header>

  <main class="container mb-5">
    <!-- Barra de Filtros com todas as Categorias -->
    <div class="card border-0 shadow-sm p-3 mb-4 bg-white">
      <div class="row align-items-center">
        <div class="col-md-5">
          <select class="form-select" id="filtroCategoria">
            <option value="todas">Todas as Categorias</option>
            <option value="Alimentação">Alimentação</option>
            <option value="Tecnologia">Tecnologia</option>
            <option value="Artesanato">Artesanato</option>
            <option value="Serviços Gerais">Serviços Gerais</option>
            <option value="Turismo">Turismo</option>
          </select>
        </div>
      </div>
    </div>

    <!-- Ponto de Montagem dos Cards Dinâmicos -->
    <div class="row" id="vitrineContainer"></div>
  </main>

  <!-- Bootstrap Bundle JS + Entrada Modular ES6 -->
  <script src="https://cdn.jsdelivr.net/npm/bootstrap@5.3.3/dist/js/bootstrap.bundle.min.js"></script>
  <script type="module" src="js/main.js"></script>
</body>
</html>
\end{boxCodigoHTML}

\begin{boxTeste}[{Checkpoint de Validação 1.1: Como Testar no Navegador}]
Abra o arquivo \texttt{index.html} no navegador usando a extensão \textbf{Live Server} do VS Code:
\begin{enumerate}[leftmargin=*]
  \item A barra verde superior com o título, o badge \textbf{"0 cadastrados"} e o botão amarelo \textbf{"Divulgar"} devem estar visíveis;
  \item O seletor de filtros com as 6 opções disponíveis (Todas + 5 categorias) deve permitir seleção;
  \item A área central estará em branco (onde os cards serão injetados pela Camada de Visão).
\end{enumerate}
\textit{Nota: Ao clicar em "Divulgar", o modal ainda não abrirá porque será estruturado no Passo 1.2.}
\end{boxTeste}

\clearpage
\begin{boxAtividade}[{Fase 1.2/5: Formulário Modal de Cadastro e Edição (index.html)}]
\textbf{Responsabilidade Arquitetural:} Fornecer a interface de diálogo acessível para inserção e atualização de serviços, contendo validações visuais nativas do HTML5 (\texttt{required}), seletores padronizados e um campo oculto (\texttt{\#servicoId}) que diferencia operações de \textit{Create} (novo) e \textit{Update} (edição).
\end{boxAtividade}

\vspace{-0.1cm}
\begin{boxCroqui}[{Croqui Visual 1.2: Anatomia do Dialogo Modal}]
+-----------------------------------------------------------------------+
|  Cadastrar Servico na Vitrine                                     [X] |
+-----------------------------------------------------------------------+
|  Nome do Empreendimento: [                                          ] |
|  Categoria: [ Selecione... v ]         Bairro: [                    ] |
|  Preco Base: [ R$ 0.00       ]         WhatsApp: [ 48999998888      ] |
|  Descricao do Servico:   [                                          ] |
|  <input type="hidden" id="servicoId">: invisivel (controle de ID)     |
+-----------------------------------------------------------------------+
|                                         [ Cancelar ]  [ Salvar ]      |
+-----------------------------------------------------------------------+
\end{boxCroqui}

\begin{boxCodigoHTML}
<!-- index.html (Parte 2: Modal de Cadastro/Edição) -->
<div class="modal fade" id="modalCadastro" tabindex="-1">
  <div class="modal-dialog modal-dialog-centered">
    <div class="modal-content shadow-lg border-0">
      <div class="modal-header">
        <h5 class="modal-title fw-bold" id="modalCadastroLabel">
          Cadastrar Serviço na Vitrine
        </h5>
        <button type="button" class="btn-close" data-bs-dismiss="modal"></button>
      </div>
      <form id="formCadastro" novalidate>
        <div class="modal-body">
          <!-- Campo oculto para ID na operação Update -->
          <input type="hidden" id="servicoId" value="">

          <div class="mb-3">
            <label for="nome" class="form-label small fw-bold">Nome do Empreendimento *</label>
            <input type="text" id="nome" class="form-control" required>
          </div>

          <div class="row g-2 mb-3">
            <div class="col-md-6">
              <label for="categoria" class="form-label small fw-bold">Categoria *</label>
              <select id="categoria" class="form-select" required>
                <option value="">Selecione...</option>
                <option value="Alimentação">Alimentação</option>
                <option value="Tecnologia">Tecnologia</option>
                <option value="Artesanato">Artesanato</option>
                <option value="Serviços Gerais">Serviços Gerais</option>
                <option value="Turismo">Turismo</option>
              </select>
            </div>
            <div class="col-md-6">
              <label for="bairro" class="form-label small fw-bold">Bairro *</label>
              <input type="text" id="bairro" class="form-control" required>
            </div>
          </div>

          <div class="row g-2 mb-3">
            <div class="col-md-6">
              <label for="precoBase" class="form-label small fw-bold">Preço Base (R$) *</label>
              <input type="number" id="precoBase" class="form-control"
                     min="0" step="0.5" required>
            </div>
            <div class="col-md-6">
              <label for="telefone" class="form-label small fw-bold">WhatsApp / Telefone *</label>
              <input type="tel" id="telefone" class="form-control"
                     placeholder="48999998888" required>
            </div>
          </div>

          <div class="mb-3">
            <label for="descricao" class="form-label small fw-bold">Descrição do Serviço *</label>
            <textarea id="descricao" class="form-control" rows="3" required></textarea>
          </div>
        </div>

        <div class="modal-footer">
          <button type="button" class="btn btn-light" data-bs-dismiss="modal">Cancelar</button>
          <button type="submit" id="btnSalvarModal" class="btn btn-success fw-bold">
            <span id="btnSalvarTexto">Salvar na Vitrine</span>
          </button>
        </div>
      </form>
    </div>
  </div>
</div>
\end{boxCodigoHTML}

\begin{boxTeste}[{Checkpoint de Validação 1.2: Como Testar a Abertura do Modal}]
Atualize o navegador e clique no botão amarelo \textbf{"Divulgar"}:
\begin{enumerate}[leftmargin=*]
  \item O modal deve abrir de forma suave e centralizada na tela;
  \item Verifique se os 6 campos possuem seus rótulos e se o select de categoria exibe todas as opções;
  \item Clique em \textbf{"Cancelar"} ou no botão \textbf{"X"} para confirmar que o fechamento nativo do Bootstrap funciona.
\end{enumerate}
\end{boxTeste}

\clearpage
\begin{boxAtividade}[{Fase 2/5: Utilitários Puros (js/utils/formatters.js)}]
\textbf{Responsabilidade Arquitetural:} Funções puras e reutilizáveis de transformação de dados. Recebem entradas brutas e retornam formatos amigáveis para o usuário brasileiro (moeda e telefone), além de realizar a sanitização de strings contra ataques XSS (\textit{Cross-Site Scripting}). Não tocam no DOM nem no armazenamento.
\end{boxAtividade}

\vspace{-0.1cm}
\begin{boxCroqui}[{Croqui Visual 2: Fluxo Puro de Entrada e Saida}]
[ "48999998888" ] --------> [ formatarTelefone() ] --------> "(48) 99999-8888"
[ 35.5          ] --------> [ formatarMoeda()    ] --------> "R$ 35,50"
[ "<script>..." ] --------> [ escapeHtml()       ] --------> "&lt;script&gt;..."
\end{boxCroqui}

\begin{boxCodigoJS}
// js/utils/formatters.js
export function formatarMoeda(valor) {
  const num = parseFloat(valor) || 0;
  return num.toLocaleString('pt-BR', {
    style: 'currency',
    currency: 'BRL'
  });
}

export function formatarTelefone(fone) {
  const digits = (fone || '').replace(/\D/g, '');
  if (digits.length === 11) {
    return digits.replace(
      /^(\d{2})(\d{5})(\d{4})/,
      '($1) $2-$3'
    );
  }
  return fone;
}

export function escapeHtml(texto) {
  const div = document.createElement('div');
  div.textContent = texto;
  return div.innerHTML;
}
\end{boxCodigoJS}

\begin{boxTeste}[{Checkpoint de Validação 2: Teste via Console do Navegador (F12)}]
Abra o navegador, pressione \textbf{F12} (ou botão direito $\rightarrow$ Inspecionar) e selecione a aba \textbf{Console}. Execute:
\begin{center}
\texttt{import('./js/utils/formatters.js').then(m => console.log(m.formatarMoeda(149.9)));}
\end{center}
\textbf{Resultado esperado no Console:} \texttt{"R\$ 149,90"}. Isso comprova que a exportação nomeada e o cálculo nativo estão 100\% operacionais!
\end{boxTeste}

\clearpage
\begin{boxAtividade}[{Fase 3/5: Camada de Serviços e Dados (js/services/vitrineService.js)}]
\textbf{Responsabilidade Arquitetural:} Isolar completamente a regra de persistência e ciclo de vida do CRUD (\textit{Create, Read, Update, Delete}). Opera estritamente sobre a chave \texttt{'garopaba\_vitrine\_servicos'} do \texttt{localStorage} do navegador, sem gerar nenhuma tag HTML nem manipular o DOM.
\end{boxAtividade}

\vspace{-0.1cm}
\begin{boxCroqui}[{Croqui Visual 3: Arquitetura da Camada de Servicos}]
[ Interface / main.js ] <--(obter, salvar, atualizar, remover)--> [ vitrineService ]
                                                                        |
                                                                        v
                                                            [ LocalStorage (Browser) ]
\end{boxCroqui}

\begin{boxCodigoJS}
// js/services/vitrineService.js
const STORAGE_KEY = 'garopaba_vitrine_servicos';

export function obterServicos() {
  const dados = localStorage.getItem(STORAGE_KEY);
  return dados ? JSON.parse(dados) : [];
}

export function salvarServico(novo) {
  const servicos = obterServicos();
  const completo = {
    id: Date.now().toString(),
    dataCadastro: new Date().toISOString(),
    ...novo
  };
  servicos.unshift(completo);
  localStorage.setItem(STORAGE_KEY, JSON.stringify(servicos));
  return completo;
}

export function atualizarServico(id, dados) {
  const servicos = obterServicos();
  const idx = servicos.findIndex(s => s.id === id);
  if (idx !== -1) {
    servicos[idx] = { ...servicos[idx], ...dados,
                      dataEdicao: new Date().toISOString() };
    localStorage.setItem(STORAGE_KEY, JSON.stringify(servicos));
    return servicos[idx];
  }
  return null;
}

export function removerServico(id) {
  const servicos = obterServicos().filter(s => s.id !== id);
  localStorage.setItem(STORAGE_KEY, JSON.stringify(servicos));
  return servicos;
}
\end{boxCodigoJS}

\begin{boxTeste}[{Checkpoint de Validação 3: Teste de Persistência no Console (F12)}]
No Console do navegador (\textbf{F12}), teste a persistência de um serviço:
\begin{center}
\texttt{import('./js/services/vitrineService.js').then(s => \{} \\
\texttt{  s.salvarServico(\{ nome: 'Oficina Garopaba', categoria: 'Serviços Gerais', precoBase: 50 \});} \\
\texttt{  console.log('Serviços no LocalStorage:', s.obterServicos());} \\
\texttt{\});}
\end{center}
\textbf{Resultado esperado:} O console imprimirá um array com o serviço recém-criado, contendo um identificador único \texttt{id} gerado automaticamente. Verifique também na aba \textbf{Application $\rightarrow$ Local Storage}.
\end{boxTeste}

\clearpage
\begin{boxAtividade}[{Fase 4/5: Camada de Visão (js/views/vitrineView.js)}]
\textbf{Responsabilidade Arquitetural:} Construir e manipular os elementos visuais na interface. Transforma objetos de dados em cards HTML (\texttt{criarCardHtml}), renderiza coleções no contêiner com tratamento de estado vazio (\texttt{renderizarCards}) e dispara mensagens de notificação flutuantes (\texttt{exibirToast}).
\end{boxAtividade}

\vspace{-0.1cm}
\begin{boxCroqui}[{Croqui Visual 4: Anatomia do Card de Servico}]
+-----------------------------------------------------------------------+
| [ Badge Categoria ]                                          R$ 35,00 | <- Topo
| Nome do Empreendimento                                                | <- Titulo
| (icone) Bairro em Garopaba                                            | <- Subtitulo
| Breve descricao dos diferenciais e produtos do negocio...             | <- Descricao
+-----------------------------------------------------------------------+
| [ WhatsApp (48) 9999-8888 ]               [ Editar ]  [ Excluir ]     | <- Acoes
+-----------------------------------------------------------------------+
\end{boxCroqui}

\begin{boxCodigoJS}
// js/views/vitrineView.js (Parte 1: Template do Card)
import { formatarMoeda, formatarTelefone, escapeHtml }
  from '../utils/formatters.js';

export function criarCardHtml(s) {
  const tel = (s.telefone || '').replace(/\D/g, '');
  return `
    <div class="col-md-6 col-lg-4 mb-4">
      <div class="card h-100 shadow-sm border-0 vitrine-card">
        <div class="card-body">
          <div class="d-flex justify-content-between mb-2">
            <span class="badge bg-success">
              ${escapeHtml(s.categoria)}
            </span>
            <span class="fw-bold text-primary">
              ${formatarMoeda(s.precoBase)}
            </span>
          </div>
          <h5 class="card-title fw-bold mb-1">${escapeHtml(s.nome)}</h5>
          <h6 class="text-muted small mb-3">
            <i class="bi bi-geo-alt-fill text-danger me-1"></i>
            ${escapeHtml(s.bairro)}
          </h6>
          <p class="card-text text-secondary small">
            ${escapeHtml(s.descricao)}
          </p>
        </div>
        <div class="card-footer bg-transparent border-0 pt-0 pb-3
                    d-flex justify-content-between align-items-center">
          <a href="https://wa.me/55${tel}" target="_blank"
             class="btn btn-sm btn-outline-success">
            <i class="bi bi-whatsapp me-1"></i>
            ${formatarTelefone(s.telefone)}
          </a>
          <div class="d-flex gap-1">
            <button class="btn btn-sm btn-outline-primary btn-editar"
                    data-id="${s.id}" title="Editar">
              <i class="bi bi-pencil-square"></i>
            </button>
            <button class="btn btn-sm btn-outline-danger btn-excluir"
                    data-id="${s.id}" title="Excluir">
              <i class="bi bi-trash"></i>
            </button>
          </div>
        </div>
      </div>
    </div>`;
}

// js/views/vitrineView.js (Parte 2: Renderização e Estado Vazio)
export function renderizarCards(servicos, container) {
  if (!servicos || servicos.length === 0) {
    container.innerHTML = `
      <div class="col-12 text-center py-5">
        <i class="bi bi-inbox fs-1 text-muted"></i>
        <h5 class="text-muted mt-2">
          Nenhum serviço cadastrado nesta categoria.
        </h5>
      </div>`;
    return;
  }
  container.innerHTML = servicos.map(criarCardHtml).join('');
}

// js/views/vitrineView.js (Parte 3: Notificações Flutuantes Toasts)
export function exibirToast(mensagem, tipo = 'success') {
  const toastEl = document.getElementById('toastNotificacao');
  const toastBody = document.getElementById('toastMensagem');
  if (toastEl && toastBody) {
    toastBody.innerText = mensagem;
    toastEl.className =
      `toast align-items-center text-bg-${tipo} border-0`;
    const toast = bootstrap.Toast.getOrCreateInstance(toastEl);
    toast.show();
  }
}
\end{boxCodigoJS}

\begin{boxTeste}[{Checkpoint de Validação 4: Teste de Renderização e Toast no Console}]
No Console do navegador (\textbf{F12}), teste o disparo da notificação:
\begin{center}
\texttt{import('./js/views/vitrineView.js').then(v => v.exibirToast('Notificação funcionando!'));}
\end{center}
\textbf{Resultado esperado:} O Toast flutuante verde surgirá no canto superior direito da tela!
\end{boxTeste}

\clearpage
\begin{boxAtividade}[{Fase 5/5: Orquestrador Central (js/main.js)}]
\textbf{Responsabilidade Arquitetural:} O controlador principal da aplicação. Importa os serviços e as visões, escuta o carregamento do DOM, conecta os eventos de submissão do formulário (\textit{Create \& Update}), a delegação de cliques nos cards (\textit{Edit \& Delete}) e a reatividade instantânea dos filtros.
\end{boxAtividade}

\vspace{-0.1cm}
\begin{boxCroqui}[{Croqui Visual 5: Ciclo de Vida da Orquestracao}]
Usuario clica/submete ---> main.js valida ---> vitrineService (localStorage)
                                                      |
Interface renderiza cards <--- vitrineView <----------+ (atualizarInterface)
\end{boxCroqui}

\begin{boxCodigoJS}
// js/main.js
import {
  obterServicos, salvarServico,
  atualizarServico, removerServico
} from './services/vitrineService.js';
import {
  renderizarCards, exibirToast
} from './views/vitrineView.js';

document.addEventListener('DOMContentLoaded', () => {
  const form = document.getElementById('formCadastro');
  const container = document.getElementById('vitrineContainer');
  const filtroCategoria = document.getElementById('filtroCategoria');
  const contadorEl = document.getElementById('totalServicosBadge');
  const modalEl = document.getElementById('modalCadastro');
  const servicoIdInput = document.getElementById('servicoId');

  function atualizarInterface() {
    const todos = obterServicos();
    const filtro = filtroCategoria.value;
    const filtrados = filtro === 'todas'
      ? todos
      : todos.filter(s => s.categoria === filtro);
    renderizarCards(filtrados, container);
    if (contadorEl) {
      contadorEl.innerText = `${todos.length} cadastrados`;
    }
  }

  // Filtros Reativos
  filtroCategoria.addEventListener('change', atualizarInterface);
  atualizarInterface();

  // Submissão do Formulário (Create & Update)
  form.addEventListener('submit', (e) => {
    e.preventDefault();
    if (!form.checkValidity()) {
      e.stopPropagation();
      form.classList.add('was-validated');
      return;
    }

    const idAtual = servicoIdInput ? servicoIdInput.value : '';
    const getVal = (id) => document.getElementById(id).value.trim();

    const dados = {
      nome: getVal('nome'),
      categoria: document.getElementById('categoria').value,
      bairro: getVal('bairro'),
      precoBase: parseFloat(getVal('precoBase')) || 0,
      telefone: getVal('telefone'),
      descricao: getVal('descricao')
    };

    if (idAtual) {
      atualizarServico(idAtual, dados);
      exibirToast('Serviço atualizado com sucesso!');
    } else {
      salvarServico(dados);
      exibirToast('Empreendimento cadastrado com sucesso!');
    }

    form.reset();
    if (servicoIdInput) servicoIdInput.value = '';
    form.classList.remove('was-validated');

    if (modalEl) {
      const modal = bootstrap.Modal.getInstance(modalEl);
      if (modal) modal.hide();
    }
    atualizarInterface();
  });

  // Delegação de Cliques nos Cards (Edit & Delete)
  container.addEventListener('click', (e) => {
    const btnEditar = e.target.closest('.btn-editar');
    if (btnEditar) {
      const id = btnEditar.getAttribute('data-id');
      const item = obterServicos().find(s => s.id === id);
      if (item) {
        if (servicoIdInput) servicoIdInput.value = item.id;
        document.getElementById('nome').value = item.nome;
        document.getElementById('categoria').value = item.categoria;
        document.getElementById('bairro').value = item.bairro;
        document.getElementById('precoBase').value = item.precoBase;
        document.getElementById('telefone').value = item.telefone;
        document.getElementById('descricao').value = item.descricao;

        if (modalEl) {
          const modal = bootstrap.Modal.getOrCreateInstance(modalEl);
          modal.show();
        }
      }
      return;
    }

    const btnExcluir = e.target.closest('.btn-excluir');
    if (btnExcluir) {
      const id = btnExcluir.getAttribute('data-id');
      if (confirm('Deseja realmente remover este serviço da vitrine?')) {
        removerServico(id);
        atualizarInterface();
        exibirToast('Serviço removido da vitrine.', 'warning');
      }
    }
  });
});
\end{boxCodigoJS}

\begin{boxTeste}[{Checkpoint de Validação 5: Teste Completo da Aplicação Ponta a Ponta}]
Parabéns! Sua aplicação está 100\% integrada. Execute o ciclo completo no navegador:
\begin{enumerate}[leftmargin=*]
  \item \textbf{Create:} Clique em "Divulgar", preencha o formulário e salve. O card deve surgir imediatamente na vitrine e o contador deve somar +1;
  \item \textbf{Filter:} Mude o filtro de categoria no select. Apenas serviços da categoria escolhida devem permanecer na tela;
  \item \textbf{Update:} Clique no ícone de lápis em um card. O modal abrirá com os dados carregados. Altere o valor e salve;
  \item \textbf{Delete:} Clique no ícone de lixeira, confirme o aviso nativo e veja o item desaparecer em tempo real!
\end{enumerate}
\end{boxTeste}

% =====================================================
% CHECKLIST DE ENTREGA DO ESTUDANTE
% =====================================================

\clearpage
\subsection{Checklist de Conclusão e Critérios Avaliativos}

Antes de realizar o commit final no repositório, certifique-se de validar todos os itens:

\begin{itemize}
\item[$\Box$] O projeto utiliza estritamente o padrão ES6 Modules com a tag:\\ \texttt{<script type="module" src="js/main.js"></script>};
\item[$\Box$] A estrutura em pastas (\texttt{js/services/}, \texttt{js/views/}, \texttt{js/utils/}) foi rigorosamente respeitada;
\item[$\Box$] O arquivo \texttt{vitrineService.js} opera puramente sobre o \texttt{localStorage} sem manipular elementos do DOM;
\item[$\Box$] As quatro operações do CRUD (\textbf{Create, Read, Update e Delete}) foram implementadas e testadas com persistência no navegador;
\item[$\Box$] O formulário exibe validações visuais nativas do Bootstrap 5 (\texttt{.is-invalid}) quando campos obrigatórios não são preenchidos;
\item[$\Box$] A vitrine reage instantaneamente aos filtros de categoria sem recarregar a página;
\item[$\Box$] O código está devidamente publicado e em pleno funcionamento no \textbf{GitHub Pages}.
\end{itemize}
